\section{Characterization of the Hyperelliptic locus}
Whenever $\tau$ is the period matrix of a hyperelliptic curve, a classical result (cf.~\cite{Th}) states that a~system of characteristics as in (\ref{eq:FS}) with the vanishing and non-vanishing properties~(\ref{vanishing}) and~(\ref{nonvanishing}) can be associated with~$\tau$ or with a conjugate of $\tau$ via the action of ${\rm Sp}(2g, \ZZ)$.

 It was not until 1984 that Mumford proved the converse statement in~\cite{mu}. We will now give a new proof of Mumford's result with a different approach that has the merit of involving Gunning's multisecant formula, as expressed in the following statement (cf.~\cite{Gu}):
 \begin{Theorem} Let $X$ be an irreducible principally polarized abelian variety of dimension~$g$, and let $A_0, \dots , A_{g+1}$ be distinct points of $X$. Suppose that $\forall\, z\in X$ the $g + 2 $ points ${\rm Th}_2(A_i + z) $ are linearly dependent.
Assume moreover the following general position condition: that there
exist some $k$ and $l$ such that for $y=(-A_k+A_l)/2$ the linear span of the points ${\rm Th}_2(A_i + y) $ for $ i = 0 , \dots , g+1$ has dimension precisely equal to $g+1$, and
not less. Then $X$ is the Jacobian of some curve $C$, and all the points $A_i$ are in the image of the Abel--Jacobi map.
\end{Theorem}

We are now in a position to prove that the above conditions are satisfied when cases~(\ref{vanishing}) and (\ref{nonvanishing}) hold; indeed, we have the following:

\begin{Theorem} Let $\tau$ be a point whose vanishing and non-vanishing theta constants are those in~\eqref{vanishing} and~\eqref{nonvanishing}; then $\tau$ is the period
matrix of a hyperelliptic Jacobian.
\end{Theorem}

\begin{proof}The period matrix turns out to be irreducible by counting the vanishing and non-vanishing theta constants or by checking the conditions of Theorem~4 in~\cite{SM}, i.e., the characters associated with the non-vanishing theta constants span the character group of
 $\Gamma_g[2,4]/\pm \Gamma_g[4,8]$.

Moreover, the vanishing conditions imply that Frobenius' formula (\ref{eq:RF})
\[
 Q[0, 0]^2 (\tau, z)=\sum_{k=0}^g Q[s_k, e_{k+1}]^2(\tau, z)
\]
 holds, hence Grushevsky's equations
 (\ref{eq:Cu1}) also hold with the coefficients being not all zero (see Proposition \ref{statement:FG}).

 We will prove that this implies the vectors ${\rm Th}_2(A_i + z) $ are linearly dependent.
 Here we take
 \[
A_0=0, \ A_1=(\tau s_0 + e_1)/2, \ \dots, \ A_k= (\tau s_{k-1} + e_k)/2, \ \dots, \ A_{g+1}= (\tau s_{g} + e_{g+1})/2,
\]
 with the corresponding characteristics being
\[
\begin{bmatrix} \ep \\ \de\end{bmatrix} =\begin{bmatrix} 0 \\ 0 \end{bmatrix} ,\begin{bmatrix} s_0 \\ e_{1}\end{bmatrix}, \dots\begin{bmatrix} s_k \\ e_{k+1} \end{bmatrix},\dots,\begin{bmatrix} s_g \\ e_{g+1} \end{bmatrix}.
\]

The $\s$-th coefficient of the vector ${\rm Th}_2(A_i + z)$ is
\[
\T[\s](\tau, z+(\tau s_{i-1} + e_i)/2)= (-1)^{\langle \s, e_i\rangle}\T[\s+s_{i-1}](\tau, z).
\]
Thus, the $2^g$ polynomials $R_\s$ in~(\ref{eq:Cu1}) give one non-trivial linear relation between the ${\rm Th}_2(A_i + z)$ for any~$z$.

We are yet to prove the general position condition. We set
\[
y=(A_1-A_0)/2= (1/4,0,\dots , 0)^{\rm t},
\]
and collect the points ${\rm Th}_2(A_i+y)$ in a $(g+2)\times 2^g$ matrix
\[
A=
 \begin{array}{ccccc}
 \begin{bmatrix} \T[0](A_0+y)&\dots&\dots& \T[\ep] (A_0+y)&\dots\\
\T[0](A_1+y)&\dots&\dots& \T[\ep] (A_1+y)&\dots\\
\dots&\dots&\dots&\dots&\dots\\
\T[0](A_{g+1}+y)&\dots&\dots& \T[\ep] (A_{g+1}+y)&\dots
 \end{bmatrix}.
 \end{array}
\]

 Now $A A^{\rm t}= B=(b_{ij}) $ with
\[
b_{ij}=\sum_{\ep} \T[\ep] (A_i+y)\T[\ep] (A_j+y)=\tt {0}{0}(A_i-A_j)\tt {0}{0}(A_i+A_j+2y).
\]

Up to constants we have
\[
\tt {0}{0}(A_0-A_j)=\tt{s_j}{e_{j+1}} \qquad {\rm and}\qquad \tt {0}{0}(A_0+A_j+2y)=\tt{s_j}{e_{j+1}+e_1}
\]
and for $1\leq i\leq j\leq g+1$
\[
\tt {0}{0}(A_i-A_j)=\tt{s_i+s_j}{e_{i+1}+e_{j+1}} \qquad {\rm and}\qquad \tt {0}{0}(A_i+A_j+2y)=\tt{s_i+s_j}{e_{i+1}+e_{j+1}+e_1}.
\]

Thus, the matrix $B$ turns into
\[ B=
 \begin{bmatrix} \tt {0}{0} \tt {0}{e_1} & \tt {0}{0} \tt {0}{e_1}&0&0&\dots&0\vspace{1mm}\\
\tt {0}{0} \tt {0}{e_1} & \tt {0}{0} \tt {0}{e_1}&0&0&\dots&0\\
0& 0&\tt {0}{0} \tt {0}{e_1}&0&\dots&0\\
\vdots & \vdots &\vdots &\ddots& &\vdots\\
0& 0&0&\dots& &\tt {0}{0} \tt {0}{e_1}
 \end{bmatrix}.
\]
Therefore, the matrix $B$ has rank $g+1$, hence $g+1\geq \operatorname{rk}(A)\geq \operatorname{rk}(B)=g+1$, which means the matrix $A$ has the required rank. Then $\tau$ is the period matrix of a Jacobian. As explained in~\cite{Gr}, since we are using points of order~2, we get that $2A_i- 2A_j=0$ as points of the Jacobian of the curve. By Abel's theorem this means there exists a function on the curve whose divisor is equal to $2A_i -2A_j$, i.e., the curve is hyperelliptic.
\end{proof}
